Biohybrid circuits between biological neurons and computational models require artificial synapses whose parameters, which are not reusable across experiments, reproduce certain dynamics. The parameterization of the graded chemical synapse model by Golowasch et al. (which separates the current into slow and fast components) is complex due to its high number of parameters, the coupling among them, and the time constraints of experiments with live neurons. This work addresses this problem through a parameterization algorithm capable of inducing in-phase or anti-phase couplings, guaranteeing realistic synaptic currents.

After reviewing synaptic models and optimization algorithms in the state of the art, Bayesian optimization is chosen for its sample efficiency and its ability to handle parameter sloppiness. Then, the evaluation function is explained, which is based on the frequency decomposition of the presynaptic potential, separating it into a slow and a fast component (\textit{spikes}), allowing for independent evaluation of each current component by combining shape similarity (with this potential) and extrema suitability. Later, when discussing Bayesian optimization, it is explained that it employs a Gaussian process with a \textit{Squared Exponential kernel} and \textit{Automatic Relevance Determination}. Furthermore, reparameterizations and logarithmic scales are incorporated into the exploration ranges to mitigate parameter sloppiness, along with dynamic bounds to adapt the exploration to each experiment.

Next, the developed validation implementations are discussed: an \textit{offline} version that operates unidirectionally with neuronal recordings which will later be evaluated by communicating a biological recording with a Hindmarsh-Rose neuronal model; and an \textit{online} module integrated into the software platform \textit{RTXI} for bidirectional communication (designed to not affect real-time operation), whose evaluation will be performed between two models. The convergence of the algorithm, the evolution of the metrics, the importance of the parameters for the target interaction, and the resulting neural dynamics will be analyzed, seeking both in-phase and anti-phase couplings.

The results demonstrate adequate convergence, inducing the desired couplings through realistic currents. In addition to the \textit{offline} version, which finishes in seconds, the \textit{online} version does so in less than 16 minutes, a duration compatible with biological experiments.
